%% =====================================================================
%%  T-13-01  Manufactured Consensus
%%  -------------------------------------------------------------------
%%  One idea per file. Core is mandatory. Slots in canonical order.
%%  Max 700 words. No layout commands. No filler. New source material
%%  for this entry goes in sources/<BOOK>/T-13-01.tex -- never by
%%  rewriting this file (docs/RULES.md R0.1).
%% =====================================================================
%% @kind: topic
%% @id: T-13-01
%% @title: Manufactured Consensus
%% @subtitle: Building the appearance of agreement without the agreement
%% @chapter: c13
%% @order: 10
%% @status: stable
%% @sources: B4
%% @updated: 2026-09-19
%% @allow-rewrite: slot migration to the seven-question format (docs/RULES.md section 2)

\topic{T-13-01}{Manufactured Consensus}{Building the appearance of agreement without the agreement}


\begin{core}
Consensus can be assembled rather than earned. Seeded audiences, planted questions, staged applause, selected testimonials and controlled comment sections all produce the appearance of general agreement, and the appearance is what does the persuading.
\end{core}

\begin{mechanism}
The construction works because agreement is normally inferred from its visible traces rather than measured directly. Nobody counts believers; people read the room. So the room is what gets built.

Three properties make the construction cheap. It needs only a minority of participants, since the rest supply themselves by copying. It needs no coordination among them, since each takes their cue from what is displayed. And it needs no persistence, since the inference is drawn at the moment of decision and the room can be different an hour later. What is required is control of what is visible at that moment --- which is why the technique concentrates on venues, formats and channels rather than on arguments.
\end{mechanism}

\begin{application}
  \item Control the venue. A format in which only agreement is visible is worth more than an argument.
  \item Seed early and visibly; the first few set the direction the rest copy.
  \item Have supporters answer dissent rather than answering it yourself --- peer correction is more persuasive and costs you nothing.
  \item Aggregate. Counts, percentages and rankings conceal their own composition.
  \item Time the display to the decision, not to the discussion.
\end{application}

\begin{feedback}
  \item Agreement is unanimous and unexplained.
  \item The supporters appear in the same place and at the same time.
  \item Dissent is absent, filtered, or answered by other audience members rather than by the presenter.
  \item The supporting voices are anonymous, aggregated or newly created.
  \item You cannot name a single person who holds the view and is not part of the presentation.
\end{feedback}

\begin{failure}
The construction is discoverable and, once discovered, discredits everything around it --- including the parts that were genuine. It also selects for the wrong kind of supporter: people willing to perform agreement are usually available to whoever pays, so a manufactured consensus is an asset that can be bought away from you. And because it substitutes for argument, it leaves the underlying case unmade, which matters the first time the audience changes.
\end{failure}

\begin{countermeasures}
  \item Leave the venue. Take the question somewhere the room was not built --- a different channel, a different audience, a private conversation.
  \item Ask for names. A consensus that will not resolve into individuals is a presentation.
  \item Look at what was excluded rather than at what was shown. Ask what would have to be filtered out for the room to look like this.
  \item Check the timestamps. A consensus that appeared all at once was built all at once.
  \item Hold your own view privately before entering, and treat any movement as something requiring an explanation.
\end{countermeasures}

\sources{B4}
\seesources
\seealso{T-15-01, T-15-02, T-23-01, T-13-02}
